\documentclass[12pt]{article}

%%% These are some packages that are useful
\usepackage{amsfonts, lipsum}
\usepackage[scr]{rsfso}
\usepackage{amsmath,amssymb, amscd,amsbsy, bbm, amsthm, enumerate}
\usepackage[top=1in, bottom=1in, left=.45in, right=.45in]{geometry}
\usepackage[unicode]{hyperref}
\usepackage{tikz, xcolor, fancyhdr}
\usepackage{graphicx, caption, subcaption}

%%% Page formatting
%\setlength{\headsep}{30pt}
\setlength{\parindent}{0pt}
\setlength{\textheight}{9in}



%%% These define our question environment and help number things correctly
\theoremstyle{definition}
\newtheorem{thm}{Theorem}
\newtheorem{question}[thm]{Question}
\newtheorem{questionS}[thm]{***Question}
\newtheorem{prop}[thm]{Proposition}
\newtheorem{lem}[thm]{Lemma}
\newtheorem{DEF}[thm]{Definition}
\newtheorem{rem}[thm]{Remark}


%%% These are some shortcuts that are handy
\def\real{{\mathbb R}}
\def\Natural{\mathbb{N}}
\def\dx{\textnormal{dx}}
\def\dy{\textnormal{dy}}
\def\dz{\textnormal{dz}}
\def\dt{\textnormal{dt}}
\def\ds{\textnormal{ds}}
\def\dw{\textnormal{dw}}
\def\Re{\textnormal{Re}}
\def\Im{\textnormal{Im}}
\def\exp{\textnormal{exp}}
\def\interior{\textnormal{interior}}
\def\al{\alpha}
\def\del{\delta}
\def\Del{\Delta}
\def\gam{\gamma}
\def\Gam{\Gamma}
\def\Om{\Omega}
\def\ep{\varepsilon}
\def\lam{\lambda}
\def\rational{{\mathbb Q}}
\def\integer{{\mathbb Z}}
\def\Q{{\mathbb Q}}
\def\Z{{\mathbb Z}}
\def\N{{\mathbb N}}
\def\R{{\mathbb R}}
\def\grad{\nabla}
\def\C{\mathcal C}
\def\P{\mathbb{P}}
\def\T{\mathcal T}
\def\I{\mathcal I}
\newcommand{\abs}[1]{\left| #1 \right|}
\newcommand{\inner}[1]{\langle #1 \rangle}
\newcommand{\norm}[1]{\left\lVert#1\right\rVert}
\newcommand{\spanvect}{\textnormal{span}}
\newcommand{\union}{\cup}
\newcommand{\Union}{\bigcup}
\def\intersect{\cap}
\def\Intersect{\bigcap}
\renewcommand{\neg}{\mathord{\sim}}














%%%%%%%%%%%%%%%%%%%%%%%%%%%%%%%%%
%%%%%%%%%%%%%%%%%%%%%%%%%%%%%%%%%
\newcommand{\name}{Duff Baker-Jarvis}
%%%%%%%%%%%%%%%%%%%%%%%%%%%%%%%%%
%%%%%%%%%%%%%%%%%%%%%%%%%%%%%%%%%
%%% Document Starts now
\begin{document}
\title{Analysis}
\date{\today}
\author{Dr. Duff Baker-Jarvis}
\maketitle

\begin{thm}[Rolle's Theorem]

\end{thm}

\begin{thm}[Mean Value Theorem]
    Let $f:\R \to \R$ be a function that is continuous on $[a,b]$ and differentiable on $(a,b)$. Then there exists a $c\in (a,b)$ such that
   $$ f'(c)= \frac{f(b)-f(a)}{b-a} $$  
\end{thm}

Proof: Consider the function $g(x)=f(x)- \frac{f(b)-f(a)}{b-a}(x-a)$. Then $ g(a)=f(a) $ and $ g(b)=f(a) $. Thus by Rolle's Theorem there exists a $c\in (a,b)$ s.t. $ g'(c)=0 $ and the rest follows. $\square$

\begin{thm}[Taylor's Theorem]
    Let $ f: \R \to \R $ be a  function that has its first $n$ derivatives continuous on $[a,b]$ and differentiable on $ (a,b) $. Fix a point to expand around, $ x_0 $. Let $ x \in [a,b]$. Then for any $n\in \N$ and any $c\in (a,b)$ there exists and $x_0\in (a,c)$ such that
    $$f(c)= \sum_{k=0}^{n} \frac{f^{(n)}(x_0)}{k!} (x-x_0)^k + \frac{f^{(n+1)}(c)}{(n+1)!}(x-x_0)^{n+1}$$  

\end{thm}

Idea of Proof:
\begin{itemize}
    \item Construct a function $F(t)$ such that 
    \begin{enumerate}
	\item $ F $ is continuous on $ (a,b) $ and differentiable on $ [a,b] $.
	\item $F(x)=f(x)$ and $ F(x_0)=f(x) $ 
	\item $F'(t)=0$ forces the correct coefficient of $ (x-x_0)^{n+1} $ 
    \end{enumerate}
\item Use the mean value theorem on $F(t)$.
\end{itemize}


Noteworthy property of Modified Taylor Series:
If $P(t)= \sum_{k=0}^{n} \frac{f^{(k)}(t)}{k!} (x-t)^k$ 
then $$ P'(t)= \frac{f^{(n+1)}(t)}{n!}(x-t)^{n} $$ 


Proof:( following Lay 28.2)
Let $ x\in [a,b]$ and fix $ n\in \N $. Let
$$F(t)= \sum_{k=0}^{n} \frac{f^{(k)}(t)}{k!}(x-t)^k + M(x-t)^{n+1}$$ 
where $M$ is the unique solution to 
$$ f(x)= \sum_{k=0}^{n} \frac{f^{(k)}(x_0)}{k!} (x-x_0)^k +M(x-x_0)^{n+1}.  $$ 
Note that by the noteworthy property mentioned above we have
$$ F'(t)= \frac{f^{(n+1)}(t)}{n!}(x-t)^n + (n+1)M(x-t)^n.$$ 
Note also that $F(x)=F(x_0)=f(x)$, so invoking the mean value theorem (or Rolle's Theorem) we obtain that there exists a $ c $ between $x$ and $ x_0 $ strictly such that
$$ F'(t)=0=\frac{f^{(n+1)}(t)}{n!}(x-t)^n + (n+1)M(x-t)^n.$$ 
Solving for $M$ gives $M= \frac{f^{(n+1)}(x_0)}{(n+1)!}$ as required. $\square$

\end{document}
